\documentclass[11pt,a4paper]{article}
\usepackage[margin=27mm]{geometry}
\usepackage[T1]{fontenc}
\usepackage{lmodern,amsmath,amssymb,amsthm,booktabs,xcolor}
\usepackage{xurl}
\usepackage[colorlinks=true,linkcolor=blue!45!black,urlcolor=blue!45!black,citecolor=blue!45!black]{hyperref}
\usepackage{microtype}
\newtheorem{theorem}{Theorem}
\newtheorem{lemma}[theorem]{Lemma}
\newtheorem{corollary}[theorem]{Corollary}
\newcommand{\pos}[1]{\left(#1\right)_+}
\newcommand{\E}{\mathbb E_M}
\newcommand{\LP}{L_{\rm raw}}
\title{Erd\H{o}s 708: global surplus and uniform upper-tail bounds}
\author{Research checkpoint}
\date{23 November 2026}
\begin{document}
\maketitle
\begin{abstract}
Maximizing surplus over all source cardinalities correctly produces a full
additive upper tail together with optimal raw dual weights for that same
tail. We prove a uniform conditional first-moment theorem for every such
prefix tail, including all prime powers and inactive factors. It gives
exponential integrability relative to the tail's own cardinality, and a
source-load truncation and price-band cleanup whose errors do not depend
on a later significant-weight cutoff. A separate actual CRT reduction
bounds the linear prefix--interval discrepancy by $O(n/\log n)$ at a
normalized positive global extremum. \textbf{The required uniform bound
$J_{2+\delta}\le C_\delta$ is not proved.} No arithmetic counterexample
to it, or proof of the unrestricted asymptotic coefficient two, is claimed.
\end{abstract}

\section{The precise objective and the optimization step}

Let $M\ge2$, let $I=\{x+1,\ldots,x+M\}$ with $x\ge0$, and put
\[
 X_w(a)=\sum_{p\le M}w_pv_p(a),\qquad 0\le w_p\le1.
\]
All valuations are raw. For $c>2$ define
\[
 J_c(w)=\sum_{a=2}^M\pos{X_w(a)-c}
          -\sum_{b\in I}\pos{X_w(b)-1}.
\]
For $A\subseteq[2,M]$, the boxed valuation-cover LP is
\[
 \LP(A,I)=\min\left\{\sum_{b\in I}z_b:
 0\le z_b\le1,\quad
 \sum_{b\in I}z_bv_p(b)\ge\sum_{a\in A}v_p(a)\right\}.
\]
The complete interval is feasible: its count of multiples of $p^j$
is at least $\lfloor M/p^j\rfloor$ at every level. LP duality and
the usual clipping of prices above one give
\begin{equation}
 \max_w J_c(w)=\max_{A\subseteq[2,M]}
                 \bigl(\LP(A,I)-c|A|\bigr).             \label{eq:global}
\end{equation}
Indeed, for each fixed $w$, maximizing the source surplus over subsets
selects exactly the strict tail $A_w=\{a:X_w(a)>c\}$. The two maxima
commute. Consequently, if $w$ maximizes $J_c$, it is a raw dual optimum
for its own full tail $A=A_w$. Otherwise a better dual for this $A$
would increase the left side of \eqref{eq:global}.
Prices on primes dividing no member of this tail may then be set to
zero: the tail is unchanged, its source loads are unchanged, and the
target hinge cannot increase. Global maximality preserves the value.
We make this normalization whenever the active support is used below.

For a paired primal optimum $z$, every positive-price row is tight,
and every target with $X_w(b)>1$ has $z_b=1$. If $n=|A|$ and
$F=\LP(A,I)=cn+J_c(w)$, global surplus maximality also gives
\begin{equation}
 \LP(A\setminus S,I)\le F-c|S|\qquad(S\subseteq A).     \label{eq:deletion}
\end{equation}
These assertions concern the same source set and interval throughout.
They do not transfer row balances to an unrelated tail.

The target requested in this investigation is a constant $C_\delta$,
independent of $M,I,w$, with $J_{2+\delta}(w)\le C_\delta$.
The raw transfer theorem in the manuscript supplied by the user
\cite{WangTransfer} already handles the fractional-to-integer step;
that step is not reproved here.

\section{Positive surplus has bounded reciprocal mean}

\begin{lemma}\label{lem:mean}
For every fixed $c>2$ there is $C(c)<\infty$ such that
$J_c(w)>0$ implies
\[
 \mu:=\sum_{p\le M}\frac{w_p}{p-1}\le C(c).
\]
This assertion does not require optimality of $w$.
\end{lemma}
\begin{proof}
Write $\E$ for averaging over $1,\ldots,M$. Prime-power counting gives
\[
 \E X=\mu-e_M(w),\qquad 0\le e_M(w)\le e_M=O(1/\log M),
\]
uniformly in the weights. To check uniformity, the floor errors cost at
most $Q(M)/M$, where $Q(M)=\#\{(p,j):p^j\le M\}$.
The omitted geometric tails cost at most $4/\sqrt M$ by splitting
at $p=\sqrt M$. Chebyshev's prime-count bound gives
$Q(M)=O(M/\log M)$.

Nonnegative divisor expansion also gives
\[
 \begin{split}
 \E X^2
 &\le \sum_p w_p^2\sum_{j\ge1}(2j-1)p^{-j}
       +2\sum_{p<q}\frac{w_pw_q}{(p-1)(q-1)}\\
 &=\mu^2+\sum_p w_p^2\frac p{(p-1)^2}
 \le\mu^2+2\mu.
 \end{split}
\]
Hence $\E(X-\mu)^2\le2\mu(1+e_M)$. Let
$A=\{a:X(a)>c\}$ and $n=|A|$. Positive surplus implies
\[
 cn<\sum_A X-\sum_I\pos{X-1}
 \le\sum_I\min(X,1)\le M,
\]
so $n/M<1/c<1/2$. For all sufficiently large $M$, $e_M\le1$.
If $\mu\ge\max(2c,33)$, then
\[
 \frac{\#\{a\le M:X(a)\le c\}}M
 \le\frac{4\mu}{(\mu-c)^2}\le\frac{16}{\mu}<\frac12,
\]
which contradicts $n/M<1/2$. The finitely many smaller values of $M$
are absorbed in $C(c)$.
\end{proof}

\section{A uniform conditional first moment on the full tail}

\begin{theorem}\label{thm:tail}
There is an absolute constant $K$ such that, for all $M$, all
$0\le w_p\le1$, and every $t\ge0$,
\begin{equation}
 \sum_{\substack{a\le M\\X(a)>t}}X(a)
 \le\bigl[t+1+K(1+\mu)\bigr]\,
       \#\{a\le M:X(a)>t\}.                            \label{eq:tailmean}
\end{equation}
In particular, no support-size, squarefreeness, prime-gap, or
source-density assumption is made.
\end{theorem}

\subsection*{A uniform rough-number mean}
For a prime $p$, set
\[
 R(T,p)=\#\{b\le T:P^-(b)\ge p\},\qquad P^-(1)=\infty.
\]
We use absolute constants in the uniform bounds
\[
 R(T,p)\gg T/\log p\quad(T\ge p^2),\qquad
 R(T,p)\ll T/\log p\quad(T\ge p).
\]
For example, Fan's uniform upper bound for $\Phi(x,y)$ and
Theorem 2.3 of \cite{Fan} give these estimates with $y=p-1$;
the finitely many small primes follow from residue-class counts.
These are uniform inequalities, not fixed-parameter asymptotics.

We claim that
\begin{equation}
 \sum_{\substack{b\le T\\P^-(b)\ge p}}X(b)
 \le K(1+\mu)R(T,p)                                  \label{eq:roughmean}
\end{equation}
with an absolute $K$. If $T<p^2$, each counted integer is one or a
prime, so its load is at most one. Otherwise expand all valuation levels:
\[
 \sum_{\substack{b\le T\\P^-(b)\ge p}}X(b)
 =\sum_{p\le q\le M}w_q\sum_{\substack{j\ge1\\q^j\le T}}
                         R(T/q^j,p).
\]
The terms with $T/q^j\ge p$ contribute at most
\[
 O\!\left(\frac T{\log p}\sum_qw_q\sum_{j\ge1}q^{-j}\right)
 =O\!\left(\frac{\mu T}{\log p}\right).
\]
For the remaining terms, $1\le T/q^j<p$, so the rough count is
exactly one. Their total is bounded by
$Q(T)=O(T/\log T)=O(T/\log p)$. Dividing by the lower bound for
$R(T,p)$ proves \eqref{eq:roughmean}.

\subsection*{Exact disjoint cells}
For every $a$ with $X(a)>t$, read \emph{all} its prime factors,
including repetitions and zero-price factors, in nondecreasing order.
Let $d$ be the first partial product with $X(d)>t$, and put
$p=P^+(d)$. Then
\[
 X(d/p)\le t<X(d)\le t+1.
\]
The numbers with this first crossing are exactly
\begin{equation}
 \{db\le M:P^-(b)\ge p\}.                             \label{eq:cell}
\end{equation}
Indeed, smaller prime factors are already included in $d$; subsequent
factors cannot change the earlier first crossing. Conversely, every
number in \eqref{eq:cell} has that crossing at $d$. Thus these cells
are pairwise disjoint and partition the entire tail.

The mean load in one cell is, by additivity and \eqref{eq:roughmean},
at most
\[
 X(d)+K(1+\mu)\le t+1+K(1+\mu).
\]
Summing over the disjoint cells proves \eqref{eq:tailmean}. No error
depending on the number of cells is incurred. Retaining zero-price
factors in the first-crossing procedure is essential for obtaining the
exact full rough-prefix cells.

\section{Exponential tails and a joint cleanup with controlled error}

Assume $\mu\le C$ and put $B=1+K(1+C)$. Define
\[
 N(t)=\#\{a\le M:X(a)>t\},\qquad
 S(t)=\sum_{a\le M}\pos{X(a)-t}.
\]
Theorem \ref{thm:tail} gives $S(t)\le BN(t)$ for every $t\ge0$.
Since $S'(t)=-N(t)$ almost everywhere,
\begin{equation}
 S(c+u)\le Bn e^{-u/B},\qquad n=N(c),\quad u\ge0.       \label{eq:exptail}
\end{equation}
In particular, for every integer $r\ge1$,
\[
 \sum_{a\in A}(X(a)-c)^r\le r!B^rn,
 \qquad
 \sum_{a\in A}e^{\lambda(X(a)-c)}\le\frac n{1-\lambda B}
 \quad(0\le\lambda<1/B).
\]
For $r\ge2$, the first estimate follows by integrating
$r(r-1)u^{r-2}S(c+u)$; the exponential estimate follows by summing
the nonnegative Taylor series. Also, for $u\ge B$,
\begin{equation}
 \sum_{X(a)>c+u}X(a)
 \le(c+u+B)n e^{1-u/B}.                               \label{eq:weightedtail}
\end{equation}
Thus for every $\varepsilon>0$ there is
$L=L(c,C,\varepsilon)$, independent of $M,n$ and the support, such that
the original weighted demand of sources with $X(a)>L$ is at most
$\varepsilon n$. Lemma \ref{lem:mean} makes the assumption
$\mu\le C(c)$ automatic at positive surplus.

Now let $w,z$ be a genuinely paired optimum for its full tail $A$.
For a set of price coordinates $E$, write
\[
 \Delta_E(a)=\sum_{p\in E}w_pv_p(a),\qquad
 D_E=\sum_{a\in A}\Delta_E(a).
\]
Fix $\gamma>0$. Delete those prices to obtain $w'$, and retain
\[
 G=\{a\in A:X_w(a)\le L,\quad X_{w'}(a)>c-\gamma\}.
\]
If $D_E\le\varepsilon\gamma n/L$, then
\begin{equation}
 \sum_{a\in A\setminus G}X_w(a)\le2\varepsilon n.      \label{eq:cleanup}
\end{equation}
The high-load sources cost at most $\varepsilon n$. Every other
discarded source loses more than $\gamma$ in the deleted band and has
load at most $L$, so these cost at most $LD_E/\gamma$.

The original paired rows give the exact same-interval identity
\[
 \sum_{b\in I}z_b\Delta_E(b)=D_E.
\]
Every new hot target remains saturated. For retained positive-price
rows, their combined oversupply relative to $G$ is
\begin{equation}
 \sum_p w'_p\left(\sum_{b\in I}z_bv_p(b)-D_p(G)\right)
 =\sum_{a\in A\setminus G}X_{w'}(a)\le2\varepsilon n.  \label{eq:rowerror}
\end{equation}
Moreover $\Phi_G(w')\ge F-D_E-2\varepsilon n$.

For any partition into $N$ price bands, one band has
$D_E\le(c+B)n/N$. The order of parameter choices is therefore
legitimate: first choose $L$ from \eqref{eq:weightedtail}, then choose
$N\ge(c+B)L/(\varepsilon\gamma)$, and only then choose the bands.
If every retained substantial price exceeds $b$, each retained source
has at most $L/b$ such occurrences, but this later bound was not used
as a multiplier in \eqref{eq:cleanup}--\eqref{eq:rowerror}.

\section{A separate actual-interval reduction of the linear discrepancy}

For completeness, there is also a genuine shift reduction; it must not
be confused with truncating valuations in the rows of a fixed interval.
For each $p\le M$ put $e_p=\lfloor\log_p M\rfloor+1$. Given a shift
$x$, preserve its residue modulo $p^{e_p}$. At most one position
$x+j$, $1\le j\le M$, is divisible by $p^{e_p}$. If it exists, choose
a lift modulo $p^{e_p+1}$ excluding the single lift that would add a
further factor $p$. CRT then gives an actual nonnegative shift $x'$
with
\[
 v_p(x'+j)=\min(v_p(x+j),e_p)\quad(p\le M,\ 1\le j\le M).
\]
Every nonnegative price vector has no larger target hinge on this new
interval. Thus a worst shift may be chosen in this finite residue
space, and weights and the paired primal may then be optimized again.

Use the established absolute bound $\LP(A,I)\le C_0|A|$ and the
ambient quadratic gate $M\ge16|A|^2\Rightarrow\LP(A,I)\le2|A|$
from the supplied manuscript \cite{WangMain}. At a positive global
surplus extremum, \eqref{eq:deletion} implies that every source prime
satisfies $p<C_0n$. To verify this, delete a source divisible by $p$.
Its cofactor bank has at least $p$ positions. If $p\ge F-c+1$, an
optimal deletion cover leaves at least one unit of residual capacity
in that bank. Purchasing this unit supplies the whole cofactor; any
remaining one-unit $p$ deficit can be repaired at cost at most one
using full-interval supply. The total is at most $F-c+2<F$, impossible.

The gate gives $M<16n^2$. With the actual capped shift and reoptimized
weights, every prime-power count differs from the prefix count by zero
or one. Hence
\begin{equation}
 \begin{split}
 0\le\sum_{b\in I}X_w(b)-\sum_{a\le M}X_w(a)
 &\le\sum_{p\le C_0n}\bigl(1+\lfloor\log_pM\rfloor\bigr)\\
 &\le2\pi(C_0n)+\pi(\sqrt M)+O(M^{1/3}\log M)\\
 &=O(n/\log n).
 \end{split}                                          \label{eq:linear}
\end{equation}
This is a linear-moment bound. It does not compare nonlinear hinges or
authorize carrying the old paired rows to the new interval.

\section{The exact limitation of this checkpoint}

The new source-tail theorem is unrestricted, and both its first-crossing
proof and the cleanup estimates have been independently checked within
this investigation. It gives a source demand budget $O_c(n)$ and
uniform integrability relative to $n$ in both sparse and dense regimes.
It does not estimate the target hinge from below by the required source
surplus. Even together with \eqref{eq:linear}, it yields at best
$J_c\le Bn$ without an additional joint arithmetic inequality; $B$ is
fixed and this error is not $o(n)$ or $O_c(1)$.

The retained source-load bound controls the number of substantial
occurrences in each source. It does not bound the global number of
different prime patterns or the number of their banks passing through
one target. The whole-row balances in \eqref{eq:rowerror} control first
weighted moments; they do not transfer source exponential moments to
targets or allocate the target hinge separately to each pattern.

Consequently the required common capacity estimate remains open in
this work. The supplied raw transfer theorem would apply after a
uniform fractional coefficient-two estimate is proved. There is no
such conclusion here, and the exact $2n$ assertion is also not claimed.

\begin{thebibliography}{9}
\bibitem{Fan} Kai Fan, \emph{Numerically explicit estimates for the
distribution of rough numbers}, Journal of Number Theory 260 (2024),
120--150. DOI: \href{https://doi.org/10.1016/j.jnt.2024.01.008}{10.1016/j.jnt.2024.01.008}.
Uniform upper bound in the introduction; uniform lower bound in Theorem 2.3.
\bibitem{WangTransfer} Wang manuscript supplied in the research request,
\emph{Raw transfer}, source file:
\url{https://github.com/science-prince/erdos-708-interval-divisor-covers/blob/main/papers/main-paper/raw-transfer.tex}.
\bibitem{WangMain} Wang manuscript supplied in the research request,
main source file:
\url{https://github.com/science-prince/erdos-708-interval-divisor-covers/blob/main/papers/main-paper/main.tex}.
\end{thebibliography}
\end{document}
